\chapter{Tuples}

\section{Why This Matters}
If a tuple is just a list you can't change, why use it?
\begin{enumerate}
    \item \textbf{Safety:} In Machine Learning, you often have data that should \textit{never} change during execution (like the shape of an image matrix: \texttt{(256, 256, 3)}). Using a tuple guarantees that a stray line of code won't accidentally mutate this critical data.
    \item \textbf{Speed and Memory:} Because tuples are immutable, Python optimizes them. They take up slightly less memory and are faster to iterate through than lists.
    \item \textbf{Dictionary Keys:} Lists cannot be used as keys in a dictionary (because they can change), but tuples can.
\end{enumerate}

\section{Explanation}
A tuple in Python is an ordered, \textbf{immutable} (unchangeable) collection of items. Tuples are created using parentheses \texttt{()} (though the parentheses are actually optional; it's the commas \texttt{,} that define a tuple). Because they are ordered, you index and slice them just like lists and strings.

\subsection{1. Syntax and Core Methods}
\begin{lstlisting}
dimensions = (1920, 1080)
hyperparameters = 0.01, 32, 64  # Parentheses are optional!

print(dimensions[0])       # 1920
# dimensions[0] = 800  <--- TypeError!
\end{lstlisting}

Because they are immutable, tuples only have \textbf{two} methods: \texttt{.count(value)} and \texttt{.index(value)}.

\subsection{2. Unpacking}
A very Pythonic feature: you can assign multiple variables at once from a tuple.
\begin{lstlisting}
model_config = ("Adam", 0.001, 100)
optimizer, learning_rate, epochs = model_config

print(learning_rate) # 0.001
\end{lstlisting}

\section{Common Mistakes}

\subsection{The Single Element Tuple Trap}
If you want to create a tuple with exactly ONE element, you \textbf{must} include a trailing comma. Parentheses alone are just treated as mathematical grouping!
\begin{lstlisting}
math_variable = (5)
print(type(math_variable))  # <class 'int'>

single_tuple = (5,)
print(type(single_tuple))   # <class 'tuple'>
\end{lstlisting}

\subsection{Nested Mutability}
A tuple itself cannot be changed. However, if a tuple \textit{contains} a mutable object (like a list), you CAN mutate the list inside the tuple!
\begin{lstlisting}
complex_tuple = (1, 2, [3, 4])
complex_tuple[2].append(5)
print(complex_tuple) # (1, 2, [3, 4, 5])
\end{lstlisting}
\textit{Why? Because the tuple is still pointing to the EXACT same list object in memory. The tuple didn't change its pointers; the list itself changed.}

\mlconnection{When you use a library like NumPy or TensorFlow, the shape of a multi-dimensional array or tensor is ALWAYS returned as a tuple (e.g., \texttt{(batch\_size, height, width, channels)}). Furthermore, many ML functions return multiple values (e.g., \texttt{return loss, accuracy}). When a function returns multiple values separated by commas, it is secretly returning a tuple!}

\section{Solutions to Exercises}

\textbf{Exercise 1: Unpacking}
\begin{lstlisting}
results = ("ResNet50", 0.92, 14.5)
name, acc, time = results
\end{lstlisting}
\textit{Solution: This unpacks the tuple into three separate variables seamlessly.}

\vspace{1em}
\textbf{Exercise 2: The Tuple Trap}
\begin{lstlisting}
lr_tuple = (0.01,)
\end{lstlisting}
\textit{Solution: Without the comma, Python evaluates \texttt{(0.01)} as a standard float.}

\vspace{1em}
\textbf{Exercise 3: Swapping Variables}
\begin{lstlisting}
x = 10
y = 20
x, y = y, x
\end{lstlisting}
\textit{Solution: This is the famous Pythonic variable swap. The right side packs a tuple \texttt{(20, 10)}, and the left side immediately unpacks it, bypassing the need for a temporary variable.}
